\documentclass{ximera}
\title{Quiz Problems and Solutions F26}

\newcommand{\pskip}{\vskip 0.1 in}

\begin{document}
\begin{abstract}
Quiz problems and solutions.
\end{abstract}
\maketitle

The answer is 24
\[
 \answer{24}
\]

The answer is 4
\[
  \answer{4}
\]

\section{Quiz 1B}

\begin{question}
Between 12:30pm and 1:00pm a beetle crawls clockwise around a circle of radius $5$ meters centered at the origin at a constant speed. It crosses the negative $y$-axis at 12:27pm, makes one complete rotation about the origin, and then crosses the line $y=x$ in the first quadrant at 12:43pm. 

{\bf Use the methods of this class} to find an expression for the function 
\[
   \theta = a(t) \, , \, 30\leq t \leq 60,
\]
that expresses the polar angle of the segment from the origin to the beetle, measured in revolutions, in terms of the number of  minutes past noon. The angle is measured postively in the counterclockwise sense from the positive $x$-axis.


\begin{itemize}
\item Explain your reasoning thoroughly and in complete sentences.

\item Draw a graph of the polar angle function to help with your explanation.

\item Include units for all numbers in all computations.

\end{itemize}

\begin{explanation}
I'll leave the graph for you, but because the beetle crawls at a constant speed, its polar angle changes at a constant rate. So the polar angle function is linear and its graph is a line \emph{segment}.  Be sure to label the axes with the appropriate variable names and units.

Both the graph and the expression for the polar angle function depend on our choice of the polar angle at 12:27pm. One possibility for the polar angle at 12:27pm is
\[
  \theta = a(27) = -1/4 \text{ rev}.
\] 
There are infinitely many other possibilities, such as $\theta=3/4$ rev. 

In any case, the change in the polar angle from 12:27pm to 12:43pm is
\[
  \Delta \theta = a(43) - a(27) = -(1+5/8) \text{ rev}= -13/8 \text{ rev}.
\]
The change is negative because the beetle turns clockwise about the origin.

So the beetle rotates about the orgin at the signed rate of 
\begin{align*}
  \omega &= \frac{\Delta \theta}{\Delta t} \\
               &= \frac{(a(43)-a(27))\text{ rev}}{(43 - 27) \text{ sec}} \\
              &=-\frac{13/8 \text{ rev}}{16 \text{ sec}} \\
              &= -\frac{13}{128} \frac{\text{rev}}{\text{sec}}.
\end{align*}

The change in its polar angle from 12:27pm to time $t$ minutes past noon is therefore
\begin{align*}
  \Delta \theta &= a(t) - a(27) \\
                     &=  \left( -\frac{13}{128} \frac{\text{rev}}{\text{sec}} \right) \left(t - 27  \right)\text{sec} \\
                     &=    -\frac{13}{128} \left(t-27 \right) \text{ rev}.      
\end{align*}

To get the polar angle at time $t$ minutes past noon, we add this change to the polar angle at 12:27pm. Omitting units from the expression, the beetle's polar angle (measured in revolutions) at time $t$ minutes past noon is therefore
\begin{align*}
 \theta &= a(t)  \\
          &= a(27) + (a(t) - a(27)) \\
          &= -\frac{1}{4} - \frac{13}{128} \left(t-27 \right) \, , \, 30 \leq t \leq 60.
\end{align*}
\end{explanation}

Find another possible expression for the polar angle function.
\begin{freeResponse}
\end{freeResponse}
\end{question}

\section{Quiz 5A}

\begin{question} \label{Qpvme}
Points $O$, $A$, $B$, $C$, and $D$ lie in a plane. Points $A$ and $B$ are $18$ meters apart. To get from $A$ to $B$ you would walk in the same direction as you would to get from $C$ to $D$.

Write an equation that express the position of $A$ relative to $O$ in terms of the vectors $\overrightarrow{CD}$ and $\overrightarrow{OB}$.

\begin{explanation}
To get from $O$ to $A$, we can go from $O$ to $B$ and then from $B$ to $A$. So 
\[
  \overrightarrow{OA} = \overrightarrow{OB} + \overrightarrow{BA}.
\]

To finish the problem, we must now express $\overrightarrow{BA}$ in terms of $\overrightarrow{CD}$.

Since to get from $A$ to $B$ you would walk in the \emph{same} direction as you would to get from $C$ to $D$, the vectors $\overrightarrow{BA}$ and $\overrightarrow{CD}$ point in \emph{opposite} directions. This means $\overrightarrow{BA}$ is a \emph{negative} scalar multiple of $\overrightarrow{CD}$.

Since $A$ and $B$ are $18$ meters apart, 
\[
  \left|  \overrightarrow{BA}\right| = 18 \text{ meters}.
\]

And since $\overrightarrow{CD}$ has length $\left| \overrightarrow{CD}\right|$ meters, we should multiply $\overrightarrow{CD}$ by the dimensionless factor
\[
       - \frac{18}{\left| \overrightarrow{CD}\right|}  
\]
to get the vector $\overrightarrow{BA}$. So
\[
     \overrightarrow{BA} =- \left( \frac{18}{\left| \overrightarrow{CD}\right|} \right) \overrightarrow{CD}
\]
and the position of $A$ relative to $O$ is given by the vector
\begin{align*}
   \overrightarrow{OA} &= \overrightarrow{OB} + \overrightarrow{BA} \\
                                 &=\overrightarrow{OB} - \left( \frac{18}{\left| \overrightarrow{CD}\right|} \right) \overrightarrow{CD}
\end{align*}

\end{explanation}
\end{question}

\section{Quiz 5B}

\begin{question} \label{Qvm93}
 A circle of radius $r$ meters rolls along the $x$-axis as a point $P$ on the circle traces the cycloid shown below. 

\begin{onlineOnly}
    \begin{center}
\desmos{lebi7kyoac}{900}{600}
\end{center}
\end{onlineOnly}

\href{https://www.desmos.com/calculator/lebi7kyoac}{Cycloid 26}



Let $\phi$ be the rotation angle of the rolling circle, measured clockwise from the upward vertrical when $P$ has coordinates $(0,2r)$.

 Use the methods of our class to write equations expressing the components of each of the following vectors in terms of $\phi$ and $r$:
\begin{enumerate}

\item The position of $C$ (the center of the rolling wheel) relative to the origin.

\item The position of $P$ relative to $C$. Start by writing an equation that expresses this vector in terms of $\phi$ and the vectors ${\bf u}$ and ${\bf v}$

\item The position of $P$ relative to $O$.
\end{enumerate}

\begin{explanation}
\begin{enumerate}
\item Since the distance $OQ$ from the origin to the point of contact of the circle and the $x$-axis is equal to the orange arclength $QR$ on the circle,
\[
  \overrightarrow{OQ} = \langle r \phi, 0 \rangle \text{ meters}.
\]
And since the vector $\overrightarrow{QC}$ has length $r$ and points in the direction of the postive $y$-axis,
\[
 \overrightarrow{QC} = \langle 0 , r \rangle \text{ meters}.
\]
So the position of $C$ relative to the origin is given by the vector (components in meters)
\begin{align*}
      \overrightarrow{OC} &= \overrightarrow{OQ} + \overrightarrow{QC} \\
                                    &= \langle r \phi, 0 \rangle + \langle 0 , r \rangle \\
                                    &= \langle r \phi , r \rangle.
\end{align*}

\item Since the angle $\phi$ from ${\bf u}$ to $\overrightarrow{CP}$ is measured in the direction from ${\bf u}$ to ${\bf v}$ (and the vectors ${\bf u}$ and ${\bf v}$ are perpendicular and have equal lengths),
\begin{align*}
     \overrightarrow{CP} &= (\cos \phi) {\bf u} + (\sin \phi) {\bf v} \\
                                   &= (\cos \phi) \langle 0, r \rangle + (\sin \phi) \langle r, 0\rangle \\
                                   &= \langle r \sin \phi , r \cos \phi \rangle ,
\end{align*} 
where all the components are measured in meters.

\item The position of $P$ relative to the origin is given by the vector (components in meters)
\begin{align*}
    \overrightarrow{OP} &=  \overrightarrow{OC} +  \overrightarrow{CP} \\
                                  & =  \langle r \phi , r \rangle + \langle r \sin \phi , r \cos \phi \rangle \\
                                  & = \langle r\phi + r \sin \phi , r + r \cos \phi \rangle.
\end{align*}
\end{enumerate}

\end{explanation}

\end{question}



\end{document}